\documentclass[activite]{thedocument}
\startdatakeys
\source{}
\cycle{}
\competencesDuSocle{}
\connaissancesRequises{}
\competencesTravaillees{}
\enddatakeys
\begin{document}\begingroup\small
    \begin{minipage}{.45\linewidth}
        \begin{center}
            \begin{tikzpicture}[scale=0.55,baseline=(O.base)]
                \coordinate(O)at(0,0);
                \draw(O)circle(4);
                \coordinate[label=above right:A](A1)at(10:4);
                \coordinate[label=above:B](A2)at(100:4);
                \coordinate[label=above left:C](A3)at(150:4);
                \foreach\ii in{1,2,3}{%
                    \node at(A\ii){$\times$};
                }
            \end{tikzpicture}
        \end{center}
    \end{minipage}%
    \begin{minipage}{.5\linewidth}
        \begin{enumerate}
            \item Appliquer le programme de tracé suivant.
            \begin{enumerate}
                \item Tracer [\textsf{AB}].
                \item Tracer [\textsf{BC}].
                \item Tracer la médiatrice (d$_1)$ de [\textsf{AB}].
                \item Tracer la médiatrice (d$_2)$ de [\textsf{BC}].
                \item Nommer \textsf{I} le point d'intersection de (d$_1)$ et (d$_2)$.
            \end{enumerate}
            \item Selon vous, que peut-on dire du point \textsf{I} ?
            \ncar2\dotedline
        \end{enumerate}
    \end{minipage}\endgroup
    \begin{flushleft}
        \textbf{3)} En s'inspirant des questions précédentes, construire le centre du
                cercle ci-dessous.
    \end{flushleft}
    \begin{center}
        \begin{tikzpicture}[scale=0.9,baseline=(O.base)]
            \coordinate(O)at(0,0);
            \draw[dotted](-5,-5)rectangle(7,5);
            \draw(O)circle(4);
        \end{tikzpicture}
    \end{center}
\end{document}